% =====================================================================
% Seção de Método -- MEST (condição L2d)
% Padrão SBC: % =====================================================================
% Seção de Método -- MEST (condição L2d)
% Padrão SBC: % =====================================================================
% Seção de Método -- MEST (condição L2d)
% Padrão SBC: % =====================================================================
% Seção de Método -- MEST (condição L2d)
% Padrão SBC: \input{secao_metodo_MEST} no corpo do artigo.
% Pacotes assumidos: amsmath, amssymb.
% =====================================================================

\section{MEST: Memória Episódica de Slots Tipados}
\label{sec:metodo}

A ideia do método cabe em uma frase: \emph{uma pergunta sobre uma conversa passada
já chega estruturada, e a memória deve ser guardada com a mesma estrutura}. Quando
alguém pergunta ``o que o James adotou em abril de 2022?'', a pergunta nomeia uma
pessoa (James), uma ação (adotar) e uma data (abril de 2022), e deixa exatamente um
campo em branco --- o objeto ---, que é justamente o que ela quer saber. Perguntas
de memória conversacional são, quase sempre, uma tupla \textbf{(quem, fez-o-quê,
com-o-quê, quando)} com um dos campos vazio. Um índice comum de palavras não
enxerga essa estrutura: ele só sabe responder ``quais trechos mencionam a palavra
$X$''. O \textbf{MEST} (\emph{Memória Episódica de Slots Tipados}) guarda cada
trecho da conversa já anotado com \emph{quem} falou, \emph{que ação} ocorreu,
\emph{sobre o quê}, de \emph{que tipo} de coisa se trata e \emph{quando} --- e
responde a pergunta procurando quem completa o campo que ficou vazio.

Em resumo, o método faz três movimentos. \textbf{(1)~Recorta} cada conversa em
\emph{episódios}: pequenos blocos de turnos consecutivos que tratam do mesmo
assunto, sendo o menor deles o par ``pergunta e resposta'' do diálogo.
\textbf{(2)~Pendura} cada episódio em cinco tipos de rótulo, os
\emph{slots}: ator, predicado, conceito, tipo e tempo. \textbf{(3)~Traduz} a pergunta para
esses mesmos cinco rótulos e pontua os episódios pelos rótulos que eles têm em
comum com a pergunta, somando ainda um canal de similaridade textual. Duas
propriedades atravessam os três movimentos: a memória é construída \emph{sem
nenhuma chamada a modelo de linguagem}, e \emph{nenhum limiar foi escolhido
olhando as respostas certas} --- todos são estimados do próprio corpus
(Seção~\ref{sec:metodo:calibracao}).

\subsection{Como a memória é construída}
\label{sec:metodo:substrato}

Uma única passagem de análise sintática por turno entrega de uma só vez quatro dos
cinco slots: sintagmas nominais (\textsf{conceito}), lemas de verbos de conteúdo
(\textsf{predicado}), menções a pessoas (\textsf{ator}) e datas (\textsf{tempo},
resolvidas contra a data da sessão e indexadas por ano, mês e dia). O quinto vem de elevar cada conceito aos seus hiperônimos na WordNet: é o
slot \textsf{tipo}, que faz ``quais \emph{comidas}'' alcançar um turno que diz
``\emph{frango assado}''.

O resultado é um grafo bipartido: de um lado os episódios, do outro os slots, e uma
aresta sempre que aquele rótulo foi \emph{literalmente pronunciado} naquele
episódio --- nada é inventado nem reescrito. Sobre esse grafo colocamos uma \emph{rotação}
(estrutura de \emph{fatgraph}): uma ordem cíclica fixa das arestas em torno de cada
vértice, $G=(V,H,\alpha,\sigma)$. A rotação é projetada, e rende dois objetos úteis. Em torno de um \textbf{slot}, a ordem é cronológica: percorrer
essa órbita é ler ``tudo o que esta pessoa, este verbo ou este tópico tocou, em
ordem'', sem precisar ranquear nada. Em torno de um \textbf{episódio}, a ordem é
sempre $\textsf{ator}\prec\textsf{predicado}\prec\textsf{conceito}\prec
\textsf{tipo}\prec\textsf{tempo}$, e portanto os pares vizinhos --- os
\emph{cantos} --- são exatamente (quem, fez-o-quê), (fez-o-quê, com-o-quê) e
(com-o-quê, quando). É a mesma tupla do primeiro parágrafo: uma pergunta é um canto
com um lado em branco.

\subsection{Como uma pergunta é respondida}
\label{sec:metodo:consulta}

A pergunta passa pelo mesmo extrator, o que produz o conjunto $\Lambda(q)$ dos
slots que ela nomeia e que existem na memória. Cada slot encontrado ``acende'' os
episódios ligados a ele, mas não com o mesmo peso: um rótulo que aparece em $200$
episódios não distingue nenhum deles, ao passo que um rótulo presente em três
praticamente resolve a busca sozinho. Por isso o peso de um slot cai com o seu
grau, e acima de um corte $h_\kappa$ ele deixa de apontar episódios e vira só um
bônus fixo --- a regra é ``\emph{um rótulo muito comum serve de filtro, nunca de
ponte}''. Somando os canais estruturais com um canal de similaridade textual, o
escore de um episódio $\varepsilon$ é
\begin{equation}
\begin{aligned}
\psi(v)&=w_{\kappa(v)}\bigl(1+\ln \deg(v)\bigr)^{-\gamma}\ \ \text{se}\ \deg(v)<h_{\kappa(v)};\quad
\psi(v)=w_{H}\ \text{caso contrário},\\
s(\varepsilon)&=w_{D}\,\delta(q,\varepsilon)+\sum_{v\in\Lambda^{*}(q,\varepsilon)}\psi(v)
+\beta\,\mathbb{1}\bigl[\varepsilon\in\mathrm{Orb}(v^{\star})\bigr],
\end{aligned}
\label{eq:psi}
\end{equation}
em que $\Lambda^{*}(q,\varepsilon)$ são os slots da pergunta (exceto os de ator)
ligados a $\varepsilon$, e $\delta(q,\varepsilon)$ é o maior cosseno entre a
pergunta e o episódio ou qualquer um de seus turnos.

O termo $\beta$ trata das perguntas que pedem uma \emph{lista} (``quais comidas
ela gosta?''). Aqui a estrutura de rotação faz algo que um buscador denso não faz:
como a órbita de um slot já \emph{é} a lista cronológica de tudo o que ele tocou,
basta tomar o slot mais raro entre os que a pergunta nomeou e promover a órbita
inteira --- sem segunda busca nem reordenação.

O ator recebe tratamento diferente dos demais, e a razão é simples: como quase toda
pergunta nomeia um dos falantes, somar um bônus por ator empurraria metade do grafo
para cima e não separaria nada. O ator, portanto, não é um canal, é um
\emph{filtro}: $\hat{s}(\varepsilon)=s(\varepsilon)\,[\pi+(1-\pi)
\min(\omega(a,\varepsilon)/\Omega,1)]$, onde $\omega(a,\varepsilon)$ mede quanto do
conteúdo daquele episódio saiu da boca de $a$. Um episódio é um par de turnos e
contém, por construção, fala dos dois participantes; então ``$a$ falou aqui'' não
diz nada, mas ``$a$ disse a maior parte disto'' diz.

Uma última assimetria é proposital: quem torna um trecho \emph{encontrável} é o
episódio, mas quem ocupa o orçamento de \emph{tokens} é o turno. Entregar episódios inteiros gastaria o
orçamento com poucos blocos grandes e perderia amplitude; entregar turnos soltos,
por outro lado, quebraria justamente o par pergunta--resposta que motivou o
episódio. A solução é pontuar turnos, deixando que cada um herde o que o seu
episódio conquistou:
\begin{equation}
s(t)=\hat{s}\bigl(\varepsilon(t)\bigr)+w_{D}\cos(q,x_t)
+\lambda\,w_{D}\!\!\max_{t'\in\varepsilon(t),\,t'\neq t}\!\!\cos(q,x_{t'}).
\label{eq:turno}
\end{equation}
O último termo é a regra do par escrita como aritmética: um turno herda uma fração
$\lambda$ da melhor semelhança entre os seus vizinhos de episódio. Assim, a
resposta ``era uma peça contemporânea chamada \emph{Finding Freedom}'' sobe junto
com o turno que menciona o assunto, mesmo não tendo palavra alguma em comum com a
pergunta. Tomam-se então os turnos em ordem decrescente até esgotar o orçamento.

\subsection{De onde vêm os limiares}
\label{sec:metodo:calibracao}

Resta uma pergunta incômoda: de onde saem $h_\kappa$, $\pi$ e os demais números?
Ajustá-los contra o conjunto anotado é prática comum, e é também o que impede levar
um método a outro corpus. Adotamos um critério simples --- \emph{o parâmetro
precisa das respostas certas para ser fixado?} --- e trocamos cada constante por um
estimador medido, durante a construção, sobre o corpus \emph{não} anotado. O corte de rótulo comum vira o quantil $0{,}99$ dos graus
daquele tipo, $h_\kappa=\lceil Q_{0,99}\{\deg(v):v\in\mathcal{S}_\kappa\}\rceil$:
uma contagem absoluta quebraria com a escala, pois num corpus dez vezes maior todo
slot a ultrapassaria e o mecanismo se desligaria sozinho. Pelo mesmo raciocínio, o
limiar de paráfrase vira um quantil alto dos cossenos realmente observados --- é
propriedade do codificador, não da tarefa --- e o filtro de ator vira
$\pi=1/|\Sigma|$ e $\Omega=\mathrm{mediana}_\varepsilon\max_a\omega(a,\varepsilon)$,
apertando-se sozinho quando há mais falantes. Já as palavras que só servem de
moldura (``conversa'', ``data'', ``tipo'') passam a ser detectadas por um contraste
de frequências, o análogo do IDF do lado da pergunta: uma palavra de tópico é
frequente nas perguntas \emph{porque} o é nas conversas, ao passo que uma de
moldura é frequente nas perguntas e ausente do que as pessoas disseram. Os pesos $w_\kappa$, $\lambda$ e $\gamma$ seguem fixos de propósito: afirmam algo
sobre o \emph{modelo} --- que um casamento de conceito informa mais do que um
palpite de hiperônimo ---, não sobre o corpus.
 no corpo do artigo.
% Pacotes assumidos: amsmath, amssymb.
% =====================================================================

\section{MEST: Memória Episódica de Slots Tipados}
\label{sec:metodo}

A ideia do método cabe em uma frase: \emph{uma pergunta sobre uma conversa passada
já chega estruturada, e a memória deve ser guardada com a mesma estrutura}. Quando
alguém pergunta ``o que o James adotou em abril de 2022?'', a pergunta nomeia uma
pessoa (James), uma ação (adotar) e uma data (abril de 2022), e deixa exatamente um
campo em branco --- o objeto ---, que é justamente o que ela quer saber. Perguntas
de memória conversacional são, quase sempre, uma tupla \textbf{(quem, fez-o-quê,
com-o-quê, quando)} com um dos campos vazio. Um índice comum de palavras não
enxerga essa estrutura: ele só sabe responder ``quais trechos mencionam a palavra
$X$''. O \textbf{MEST} (\emph{Memória Episódica de Slots Tipados}) guarda cada
trecho da conversa já anotado com \emph{quem} falou, \emph{que ação} ocorreu,
\emph{sobre o quê}, de \emph{que tipo} de coisa se trata e \emph{quando} --- e
responde a pergunta procurando quem completa o campo que ficou vazio.

Em resumo, o método faz três movimentos. \textbf{(1)~Recorta} cada conversa em
\emph{episódios}: pequenos blocos de turnos consecutivos que tratam do mesmo
assunto, sendo o menor deles o par ``pergunta e resposta'' do diálogo.
\textbf{(2)~Pendura} cada episódio em cinco tipos de rótulo, os
\emph{slots}: ator, predicado, conceito, tipo e tempo. \textbf{(3)~Traduz} a pergunta para
esses mesmos cinco rótulos e pontua os episódios pelos rótulos que eles têm em
comum com a pergunta, somando ainda um canal de similaridade textual. Duas
propriedades atravessam os três movimentos: a memória é construída \emph{sem
nenhuma chamada a modelo de linguagem}, e \emph{nenhum limiar foi escolhido
olhando as respostas certas} --- todos são estimados do próprio corpus
(Seção~\ref{sec:metodo:calibracao}).

\subsection{Como a memória é construída}
\label{sec:metodo:substrato}

Uma única passagem de análise sintática por turno entrega de uma só vez quatro dos
cinco slots: sintagmas nominais (\textsf{conceito}), lemas de verbos de conteúdo
(\textsf{predicado}), menções a pessoas (\textsf{ator}) e datas (\textsf{tempo},
resolvidas contra a data da sessão e indexadas por ano, mês e dia). O quinto vem de elevar cada conceito aos seus hiperônimos na WordNet: é o
slot \textsf{tipo}, que faz ``quais \emph{comidas}'' alcançar um turno que diz
``\emph{frango assado}''.

O resultado é um grafo bipartido: de um lado os episódios, do outro os slots, e uma
aresta sempre que aquele rótulo foi \emph{literalmente pronunciado} naquele
episódio --- nada é inventado nem reescrito. Sobre esse grafo colocamos uma \emph{rotação}
(estrutura de \emph{fatgraph}): uma ordem cíclica fixa das arestas em torno de cada
vértice, $G=(V,H,\alpha,\sigma)$. A rotação é projetada, e rende dois objetos úteis. Em torno de um \textbf{slot}, a ordem é cronológica: percorrer
essa órbita é ler ``tudo o que esta pessoa, este verbo ou este tópico tocou, em
ordem'', sem precisar ranquear nada. Em torno de um \textbf{episódio}, a ordem é
sempre $\textsf{ator}\prec\textsf{predicado}\prec\textsf{conceito}\prec
\textsf{tipo}\prec\textsf{tempo}$, e portanto os pares vizinhos --- os
\emph{cantos} --- são exatamente (quem, fez-o-quê), (fez-o-quê, com-o-quê) e
(com-o-quê, quando). É a mesma tupla do primeiro parágrafo: uma pergunta é um canto
com um lado em branco.

\subsection{Como uma pergunta é respondida}
\label{sec:metodo:consulta}

A pergunta passa pelo mesmo extrator, o que produz o conjunto $\Lambda(q)$ dos
slots que ela nomeia e que existem na memória. Cada slot encontrado ``acende'' os
episódios ligados a ele, mas não com o mesmo peso: um rótulo que aparece em $200$
episódios não distingue nenhum deles, ao passo que um rótulo presente em três
praticamente resolve a busca sozinho. Por isso o peso de um slot cai com o seu
grau, e acima de um corte $h_\kappa$ ele deixa de apontar episódios e vira só um
bônus fixo --- a regra é ``\emph{um rótulo muito comum serve de filtro, nunca de
ponte}''. Somando os canais estruturais com um canal de similaridade textual, o
escore de um episódio $\varepsilon$ é
\begin{equation}
\begin{aligned}
\psi(v)&=w_{\kappa(v)}\bigl(1+\ln \deg(v)\bigr)^{-\gamma}\ \ \text{se}\ \deg(v)<h_{\kappa(v)};\quad
\psi(v)=w_{H}\ \text{caso contrário},\\
s(\varepsilon)&=w_{D}\,\delta(q,\varepsilon)+\sum_{v\in\Lambda^{*}(q,\varepsilon)}\psi(v)
+\beta\,\mathbb{1}\bigl[\varepsilon\in\mathrm{Orb}(v^{\star})\bigr],
\end{aligned}
\label{eq:psi}
\end{equation}
em que $\Lambda^{*}(q,\varepsilon)$ são os slots da pergunta (exceto os de ator)
ligados a $\varepsilon$, e $\delta(q,\varepsilon)$ é o maior cosseno entre a
pergunta e o episódio ou qualquer um de seus turnos.

O termo $\beta$ trata das perguntas que pedem uma \emph{lista} (``quais comidas
ela gosta?''). Aqui a estrutura de rotação faz algo que um buscador denso não faz:
como a órbita de um slot já \emph{é} a lista cronológica de tudo o que ele tocou,
basta tomar o slot mais raro entre os que a pergunta nomeou e promover a órbita
inteira --- sem segunda busca nem reordenação.

O ator recebe tratamento diferente dos demais, e a razão é simples: como quase toda
pergunta nomeia um dos falantes, somar um bônus por ator empurraria metade do grafo
para cima e não separaria nada. O ator, portanto, não é um canal, é um
\emph{filtro}: $\hat{s}(\varepsilon)=s(\varepsilon)\,[\pi+(1-\pi)
\min(\omega(a,\varepsilon)/\Omega,1)]$, onde $\omega(a,\varepsilon)$ mede quanto do
conteúdo daquele episódio saiu da boca de $a$. Um episódio é um par de turnos e
contém, por construção, fala dos dois participantes; então ``$a$ falou aqui'' não
diz nada, mas ``$a$ disse a maior parte disto'' diz.

Uma última assimetria é proposital: quem torna um trecho \emph{encontrável} é o
episódio, mas quem ocupa o orçamento de \emph{tokens} é o turno. Entregar episódios inteiros gastaria o
orçamento com poucos blocos grandes e perderia amplitude; entregar turnos soltos,
por outro lado, quebraria justamente o par pergunta--resposta que motivou o
episódio. A solução é pontuar turnos, deixando que cada um herde o que o seu
episódio conquistou:
\begin{equation}
s(t)=\hat{s}\bigl(\varepsilon(t)\bigr)+w_{D}\cos(q,x_t)
+\lambda\,w_{D}\!\!\max_{t'\in\varepsilon(t),\,t'\neq t}\!\!\cos(q,x_{t'}).
\label{eq:turno}
\end{equation}
O último termo é a regra do par escrita como aritmética: um turno herda uma fração
$\lambda$ da melhor semelhança entre os seus vizinhos de episódio. Assim, a
resposta ``era uma peça contemporânea chamada \emph{Finding Freedom}'' sobe junto
com o turno que menciona o assunto, mesmo não tendo palavra alguma em comum com a
pergunta. Tomam-se então os turnos em ordem decrescente até esgotar o orçamento.

\subsection{De onde vêm os limiares}
\label{sec:metodo:calibracao}

Resta uma pergunta incômoda: de onde saem $h_\kappa$, $\pi$ e os demais números?
Ajustá-los contra o conjunto anotado é prática comum, e é também o que impede levar
um método a outro corpus. Adotamos um critério simples --- \emph{o parâmetro
precisa das respostas certas para ser fixado?} --- e trocamos cada constante por um
estimador medido, durante a construção, sobre o corpus \emph{não} anotado. O corte de rótulo comum vira o quantil $0{,}99$ dos graus
daquele tipo, $h_\kappa=\lceil Q_{0,99}\{\deg(v):v\in\mathcal{S}_\kappa\}\rceil$:
uma contagem absoluta quebraria com a escala, pois num corpus dez vezes maior todo
slot a ultrapassaria e o mecanismo se desligaria sozinho. Pelo mesmo raciocínio, o
limiar de paráfrase vira um quantil alto dos cossenos realmente observados --- é
propriedade do codificador, não da tarefa --- e o filtro de ator vira
$\pi=1/|\Sigma|$ e $\Omega=\mathrm{mediana}_\varepsilon\max_a\omega(a,\varepsilon)$,
apertando-se sozinho quando há mais falantes. Já as palavras que só servem de
moldura (``conversa'', ``data'', ``tipo'') passam a ser detectadas por um contraste
de frequências, o análogo do IDF do lado da pergunta: uma palavra de tópico é
frequente nas perguntas \emph{porque} o é nas conversas, ao passo que uma de
moldura é frequente nas perguntas e ausente do que as pessoas disseram. Os pesos $w_\kappa$, $\lambda$ e $\gamma$ seguem fixos de propósito: afirmam algo
sobre o \emph{modelo} --- que um casamento de conceito informa mais do que um
palpite de hiperônimo ---, não sobre o corpus.
 no corpo do artigo.
% Pacotes assumidos: amsmath, amssymb.
% =====================================================================

\section{MEST: Memória Episódica de Slots Tipados}
\label{sec:metodo}

A ideia do método cabe em uma frase: \emph{uma pergunta sobre uma conversa passada
já chega estruturada, e a memória deve ser guardada com a mesma estrutura}. Quando
alguém pergunta ``o que o James adotou em abril de 2022?'', a pergunta nomeia uma
pessoa (James), uma ação (adotar) e uma data (abril de 2022), e deixa exatamente um
campo em branco --- o objeto ---, que é justamente o que ela quer saber. Perguntas
de memória conversacional são, quase sempre, uma tupla \textbf{(quem, fez-o-quê,
com-o-quê, quando)} com um dos campos vazio. Um índice comum de palavras não
enxerga essa estrutura: ele só sabe responder ``quais trechos mencionam a palavra
$X$''. O \textbf{MEST} (\emph{Memória Episódica de Slots Tipados}) guarda cada
trecho da conversa já anotado com \emph{quem} falou, \emph{que ação} ocorreu,
\emph{sobre o quê}, de \emph{que tipo} de coisa se trata e \emph{quando} --- e
responde a pergunta procurando quem completa o campo que ficou vazio.

Em resumo, o método faz três movimentos. \textbf{(1)~Recorta} cada conversa em
\emph{episódios}: pequenos blocos de turnos consecutivos que tratam do mesmo
assunto, sendo o menor deles o par ``pergunta e resposta'' do diálogo.
\textbf{(2)~Pendura} cada episódio em cinco tipos de rótulo, os
\emph{slots}: ator, predicado, conceito, tipo e tempo. \textbf{(3)~Traduz} a pergunta para
esses mesmos cinco rótulos e pontua os episódios pelos rótulos que eles têm em
comum com a pergunta, somando ainda um canal de similaridade textual. Duas
propriedades atravessam os três movimentos: a memória é construída \emph{sem
nenhuma chamada a modelo de linguagem}, e \emph{nenhum limiar foi escolhido
olhando as respostas certas} --- todos são estimados do próprio corpus
(Seção~\ref{sec:metodo:calibracao}).

\subsection{Como a memória é construída}
\label{sec:metodo:substrato}

Uma única passagem de análise sintática por turno entrega de uma só vez quatro dos
cinco slots: sintagmas nominais (\textsf{conceito}), lemas de verbos de conteúdo
(\textsf{predicado}), menções a pessoas (\textsf{ator}) e datas (\textsf{tempo},
resolvidas contra a data da sessão e indexadas por ano, mês e dia). O quinto vem de elevar cada conceito aos seus hiperônimos na WordNet: é o
slot \textsf{tipo}, que faz ``quais \emph{comidas}'' alcançar um turno que diz
``\emph{frango assado}''.

O resultado é um grafo bipartido: de um lado os episódios, do outro os slots, e uma
aresta sempre que aquele rótulo foi \emph{literalmente pronunciado} naquele
episódio --- nada é inventado nem reescrito. Sobre esse grafo colocamos uma \emph{rotação}
(estrutura de \emph{fatgraph}): uma ordem cíclica fixa das arestas em torno de cada
vértice, $G=(V,H,\alpha,\sigma)$. A rotação é projetada, e rende dois objetos úteis. Em torno de um \textbf{slot}, a ordem é cronológica: percorrer
essa órbita é ler ``tudo o que esta pessoa, este verbo ou este tópico tocou, em
ordem'', sem precisar ranquear nada. Em torno de um \textbf{episódio}, a ordem é
sempre $\textsf{ator}\prec\textsf{predicado}\prec\textsf{conceito}\prec
\textsf{tipo}\prec\textsf{tempo}$, e portanto os pares vizinhos --- os
\emph{cantos} --- são exatamente (quem, fez-o-quê), (fez-o-quê, com-o-quê) e
(com-o-quê, quando). É a mesma tupla do primeiro parágrafo: uma pergunta é um canto
com um lado em branco.

\subsection{Como uma pergunta é respondida}
\label{sec:metodo:consulta}

A pergunta passa pelo mesmo extrator, o que produz o conjunto $\Lambda(q)$ dos
slots que ela nomeia e que existem na memória. Cada slot encontrado ``acende'' os
episódios ligados a ele, mas não com o mesmo peso: um rótulo que aparece em $200$
episódios não distingue nenhum deles, ao passo que um rótulo presente em três
praticamente resolve a busca sozinho. Por isso o peso de um slot cai com o seu
grau, e acima de um corte $h_\kappa$ ele deixa de apontar episódios e vira só um
bônus fixo --- a regra é ``\emph{um rótulo muito comum serve de filtro, nunca de
ponte}''. Somando os canais estruturais com um canal de similaridade textual, o
escore de um episódio $\varepsilon$ é
\begin{equation}
\begin{aligned}
\psi(v)&=w_{\kappa(v)}\bigl(1+\ln \deg(v)\bigr)^{-\gamma}\ \ \text{se}\ \deg(v)<h_{\kappa(v)};\quad
\psi(v)=w_{H}\ \text{caso contrário},\\
s(\varepsilon)&=w_{D}\,\delta(q,\varepsilon)+\sum_{v\in\Lambda^{*}(q,\varepsilon)}\psi(v)
+\beta\,\mathbb{1}\bigl[\varepsilon\in\mathrm{Orb}(v^{\star})\bigr],
\end{aligned}
\label{eq:psi}
\end{equation}
em que $\Lambda^{*}(q,\varepsilon)$ são os slots da pergunta (exceto os de ator)
ligados a $\varepsilon$, e $\delta(q,\varepsilon)$ é o maior cosseno entre a
pergunta e o episódio ou qualquer um de seus turnos.

O termo $\beta$ trata das perguntas que pedem uma \emph{lista} (``quais comidas
ela gosta?''). Aqui a estrutura de rotação faz algo que um buscador denso não faz:
como a órbita de um slot já \emph{é} a lista cronológica de tudo o que ele tocou,
basta tomar o slot mais raro entre os que a pergunta nomeou e promover a órbita
inteira --- sem segunda busca nem reordenação.

O ator recebe tratamento diferente dos demais, e a razão é simples: como quase toda
pergunta nomeia um dos falantes, somar um bônus por ator empurraria metade do grafo
para cima e não separaria nada. O ator, portanto, não é um canal, é um
\emph{filtro}: $\hat{s}(\varepsilon)=s(\varepsilon)\,[\pi+(1-\pi)
\min(\omega(a,\varepsilon)/\Omega,1)]$, onde $\omega(a,\varepsilon)$ mede quanto do
conteúdo daquele episódio saiu da boca de $a$. Um episódio é um par de turnos e
contém, por construção, fala dos dois participantes; então ``$a$ falou aqui'' não
diz nada, mas ``$a$ disse a maior parte disto'' diz.

Uma última assimetria é proposital: quem torna um trecho \emph{encontrável} é o
episódio, mas quem ocupa o orçamento de \emph{tokens} é o turno. Entregar episódios inteiros gastaria o
orçamento com poucos blocos grandes e perderia amplitude; entregar turnos soltos,
por outro lado, quebraria justamente o par pergunta--resposta que motivou o
episódio. A solução é pontuar turnos, deixando que cada um herde o que o seu
episódio conquistou:
\begin{equation}
s(t)=\hat{s}\bigl(\varepsilon(t)\bigr)+w_{D}\cos(q,x_t)
+\lambda\,w_{D}\!\!\max_{t'\in\varepsilon(t),\,t'\neq t}\!\!\cos(q,x_{t'}).
\label{eq:turno}
\end{equation}
O último termo é a regra do par escrita como aritmética: um turno herda uma fração
$\lambda$ da melhor semelhança entre os seus vizinhos de episódio. Assim, a
resposta ``era uma peça contemporânea chamada \emph{Finding Freedom}'' sobe junto
com o turno que menciona o assunto, mesmo não tendo palavra alguma em comum com a
pergunta. Tomam-se então os turnos em ordem decrescente até esgotar o orçamento.

\subsection{De onde vêm os limiares}
\label{sec:metodo:calibracao}

Resta uma pergunta incômoda: de onde saem $h_\kappa$, $\pi$ e os demais números?
Ajustá-los contra o conjunto anotado é prática comum, e é também o que impede levar
um método a outro corpus. Adotamos um critério simples --- \emph{o parâmetro
precisa das respostas certas para ser fixado?} --- e trocamos cada constante por um
estimador medido, durante a construção, sobre o corpus \emph{não} anotado. O corte de rótulo comum vira o quantil $0{,}99$ dos graus
daquele tipo, $h_\kappa=\lceil Q_{0,99}\{\deg(v):v\in\mathcal{S}_\kappa\}\rceil$:
uma contagem absoluta quebraria com a escala, pois num corpus dez vezes maior todo
slot a ultrapassaria e o mecanismo se desligaria sozinho. Pelo mesmo raciocínio, o
limiar de paráfrase vira um quantil alto dos cossenos realmente observados --- é
propriedade do codificador, não da tarefa --- e o filtro de ator vira
$\pi=1/|\Sigma|$ e $\Omega=\mathrm{mediana}_\varepsilon\max_a\omega(a,\varepsilon)$,
apertando-se sozinho quando há mais falantes. Já as palavras que só servem de
moldura (``conversa'', ``data'', ``tipo'') passam a ser detectadas por um contraste
de frequências, o análogo do IDF do lado da pergunta: uma palavra de tópico é
frequente nas perguntas \emph{porque} o é nas conversas, ao passo que uma de
moldura é frequente nas perguntas e ausente do que as pessoas disseram. Os pesos $w_\kappa$, $\lambda$ e $\gamma$ seguem fixos de propósito: afirmam algo
sobre o \emph{modelo} --- que um casamento de conceito informa mais do que um
palpite de hiperônimo ---, não sobre o corpus.
 no corpo do artigo.
% Pacotes assumidos: amsmath, amssymb.
% =====================================================================

\section{MEST: Memória Episódica de Slots Tipados}
\label{sec:metodo}

A ideia do método cabe em uma frase: \emph{uma pergunta sobre uma conversa passada
já chega estruturada, e a memória deve ser guardada com a mesma estrutura}. Quando
alguém pergunta ``o que o James adotou em abril de 2022?'', a pergunta nomeia uma
pessoa (James), uma ação (adotar) e uma data (abril de 2022), e deixa exatamente um
campo em branco --- o objeto ---, que é justamente o que ela quer saber. Perguntas
de memória conversacional são, quase sempre, uma tupla \textbf{(quem, fez-o-quê,
com-o-quê, quando)} com um dos campos vazio. Um índice comum de palavras não
enxerga essa estrutura: ele só sabe responder ``quais trechos mencionam a palavra
$X$''. O \textbf{MEST} (\emph{Memória Episódica de Slots Tipados}) guarda cada
trecho da conversa já anotado com \emph{quem} falou, \emph{que ação} ocorreu,
\emph{sobre o quê}, de \emph{que tipo} de coisa se trata e \emph{quando} --- e
responde a pergunta procurando quem completa o campo que ficou vazio.

Em resumo, o método faz três movimentos. \textbf{(1)~Recorta} cada conversa em
\emph{episódios}: pequenos blocos de turnos consecutivos que tratam do mesmo
assunto, sendo o menor deles o par ``pergunta e resposta'' do diálogo.
\textbf{(2)~Pendura} cada episódio em cinco tipos de rótulo, os
\emph{slots}: ator, predicado, conceito, tipo e tempo. \textbf{(3)~Traduz} a pergunta para
esses mesmos cinco rótulos e pontua os episódios pelos rótulos que eles têm em
comum com a pergunta, somando ainda um canal de similaridade textual. Duas
propriedades atravessam os três movimentos: a memória é construída \emph{sem
nenhuma chamada a modelo de linguagem}, e \emph{nenhum limiar foi escolhido
olhando as respostas certas} --- todos são estimados do próprio corpus
(Seção~\ref{sec:metodo:calibracao}).

\subsection{Como a memória é construída}
\label{sec:metodo:substrato}

Uma única passagem de análise sintática por turno entrega de uma só vez quatro dos
cinco slots: sintagmas nominais (\textsf{conceito}), lemas de verbos de conteúdo
(\textsf{predicado}), menções a pessoas (\textsf{ator}) e datas (\textsf{tempo},
resolvidas contra a data da sessão e indexadas por ano, mês e dia). O quinto vem de elevar cada conceito aos seus hiperônimos na WordNet: é o
slot \textsf{tipo}, que faz ``quais \emph{comidas}'' alcançar um turno que diz
``\emph{frango assado}''.

O resultado é um grafo bipartido: de um lado os episódios, do outro os slots, e uma
aresta sempre que aquele rótulo foi \emph{literalmente pronunciado} naquele
episódio --- nada é inventado nem reescrito. Sobre esse grafo colocamos uma \emph{rotação}
(estrutura de \emph{fatgraph}): uma ordem cíclica fixa das arestas em torno de cada
vértice, $G=(V,H,\alpha,\sigma)$. A rotação é projetada, e rende dois objetos úteis. Em torno de um \textbf{slot}, a ordem é cronológica: percorrer
essa órbita é ler ``tudo o que esta pessoa, este verbo ou este tópico tocou, em
ordem'', sem precisar ranquear nada. Em torno de um \textbf{episódio}, a ordem é
sempre $\textsf{ator}\prec\textsf{predicado}\prec\textsf{conceito}\prec
\textsf{tipo}\prec\textsf{tempo}$, e portanto os pares vizinhos --- os
\emph{cantos} --- são exatamente (quem, fez-o-quê), (fez-o-quê, com-o-quê) e
(com-o-quê, quando). É a mesma tupla do primeiro parágrafo: uma pergunta é um canto
com um lado em branco.

\subsection{Como uma pergunta é respondida}
\label{sec:metodo:consulta}

A pergunta passa pelo mesmo extrator, o que produz o conjunto $\Lambda(q)$ dos
slots que ela nomeia e que existem na memória. Cada slot encontrado ``acende'' os
episódios ligados a ele, mas não com o mesmo peso: um rótulo que aparece em $200$
episódios não distingue nenhum deles, ao passo que um rótulo presente em três
praticamente resolve a busca sozinho. Por isso o peso de um slot cai com o seu
grau, e acima de um corte $h_\kappa$ ele deixa de apontar episódios e vira só um
bônus fixo --- a regra é ``\emph{um rótulo muito comum serve de filtro, nunca de
ponte}''. Somando os canais estruturais com um canal de similaridade textual, o
escore de um episódio $\varepsilon$ é
\begin{equation}
\begin{aligned}
\psi(v)&=w_{\kappa(v)}\bigl(1+\ln \deg(v)\bigr)^{-\gamma}\ \ \text{se}\ \deg(v)<h_{\kappa(v)};\quad
\psi(v)=w_{H}\ \text{caso contrário},\\
s(\varepsilon)&=w_{D}\,\delta(q,\varepsilon)+\sum_{v\in\Lambda^{*}(q,\varepsilon)}\psi(v)
+\beta\,\mathbb{1}\bigl[\varepsilon\in\mathrm{Orb}(v^{\star})\bigr],
\end{aligned}
\label{eq:psi}
\end{equation}
em que $\Lambda^{*}(q,\varepsilon)$ são os slots da pergunta (exceto os de ator)
ligados a $\varepsilon$, e $\delta(q,\varepsilon)$ é o maior cosseno entre a
pergunta e o episódio ou qualquer um de seus turnos.

O termo $\beta$ trata das perguntas que pedem uma \emph{lista} (``quais comidas
ela gosta?''). Aqui a estrutura de rotação faz algo que um buscador denso não faz:
como a órbita de um slot já \emph{é} a lista cronológica de tudo o que ele tocou,
basta tomar o slot mais raro entre os que a pergunta nomeou e promover a órbita
inteira --- sem segunda busca nem reordenação.

O ator recebe tratamento diferente dos demais, e a razão é simples: como quase toda
pergunta nomeia um dos falantes, somar um bônus por ator empurraria metade do grafo
para cima e não separaria nada. O ator, portanto, não é um canal, é um
\emph{filtro}: $\hat{s}(\varepsilon)=s(\varepsilon)\,[\pi+(1-\pi)
\min(\omega(a,\varepsilon)/\Omega,1)]$, onde $\omega(a,\varepsilon)$ mede quanto do
conteúdo daquele episódio saiu da boca de $a$. Um episódio é um par de turnos e
contém, por construção, fala dos dois participantes; então ``$a$ falou aqui'' não
diz nada, mas ``$a$ disse a maior parte disto'' diz.

Uma última assimetria é proposital: quem torna um trecho \emph{encontrável} é o
episódio, mas quem ocupa o orçamento de \emph{tokens} é o turno. Entregar episódios inteiros gastaria o
orçamento com poucos blocos grandes e perderia amplitude; entregar turnos soltos,
por outro lado, quebraria justamente o par pergunta--resposta que motivou o
episódio. A solução é pontuar turnos, deixando que cada um herde o que o seu
episódio conquistou:
\begin{equation}
s(t)=\hat{s}\bigl(\varepsilon(t)\bigr)+w_{D}\cos(q,x_t)
+\lambda\,w_{D}\!\!\max_{t'\in\varepsilon(t),\,t'\neq t}\!\!\cos(q,x_{t'}).
\label{eq:turno}
\end{equation}
O último termo é a regra do par escrita como aritmética: um turno herda uma fração
$\lambda$ da melhor semelhança entre os seus vizinhos de episódio. Assim, a
resposta ``era uma peça contemporânea chamada \emph{Finding Freedom}'' sobe junto
com o turno que menciona o assunto, mesmo não tendo palavra alguma em comum com a
pergunta. Tomam-se então os turnos em ordem decrescente até esgotar o orçamento.

\subsection{De onde vêm os limiares}
\label{sec:metodo:calibracao}

Resta uma pergunta incômoda: de onde saem $h_\kappa$, $\pi$ e os demais números?
Ajustá-los contra o conjunto anotado é prática comum, e é também o que impede levar
um método a outro corpus. Adotamos um critério simples --- \emph{o parâmetro
precisa das respostas certas para ser fixado?} --- e trocamos cada constante por um
estimador medido, durante a construção, sobre o corpus \emph{não} anotado. O corte de rótulo comum vira o quantil $0{,}99$ dos graus
daquele tipo, $h_\kappa=\lceil Q_{0,99}\{\deg(v):v\in\mathcal{S}_\kappa\}\rceil$:
uma contagem absoluta quebraria com a escala, pois num corpus dez vezes maior todo
slot a ultrapassaria e o mecanismo se desligaria sozinho. Pelo mesmo raciocínio, o
limiar de paráfrase vira um quantil alto dos cossenos realmente observados --- é
propriedade do codificador, não da tarefa --- e o filtro de ator vira
$\pi=1/|\Sigma|$ e $\Omega=\mathrm{mediana}_\varepsilon\max_a\omega(a,\varepsilon)$,
apertando-se sozinho quando há mais falantes. Já as palavras que só servem de
moldura (``conversa'', ``data'', ``tipo'') passam a ser detectadas por um contraste
de frequências, o análogo do IDF do lado da pergunta: uma palavra de tópico é
frequente nas perguntas \emph{porque} o é nas conversas, ao passo que uma de
moldura é frequente nas perguntas e ausente do que as pessoas disseram. Os pesos $w_\kappa$, $\lambda$ e $\gamma$ seguem fixos de propósito: afirmam algo
sobre o \emph{modelo} --- que um casamento de conceito informa mais do que um
palpite de hiperônimo ---, não sobre o corpus.
